\documentclass[11pt]{article}
\usepackage{cmap}
\usepackage[utf8]{inputenc}
\usepackage[T1]{fontenc}
\usepackage[spanish]{babel}
\usepackage[margin=2.7cm]{geometry}
\usepackage{graphicx}
\usepackage{tikz}
\usetikzlibrary{arrows.meta,positioning,fit,backgrounds}
\usepackage{booktabs}
\usepackage{longtable}
\usepackage{xcolor}
\usepackage{hyperref}
\usepackage{enumitem}
\usepackage{amsfonts}

\hypersetup{
    colorlinks=true,
    linkcolor=blue!50!black,
    urlcolor=blue!50!black,
    citecolor=blue!50!black,
    pdftitle={FinHive: arquitectura multiagente jerarquica financiera sobre Databricks},
    pdfauthor={Matias Fabian Adell}
}

\definecolor{boxblue}{RGB}{235,242,252}
\definecolor{boxborder}{RGB}{70,110,180}
\definecolor{warnbox}{RGB}{253,240,230}
\definecolor{warnborder}{RGB}{200,110,40}

\newcommand{\hlr}[1]{\textcolor{red!70!black}{#1}}
\newcommand{\hlb}[1]{\textcolor{blue!50!black}{#1}}
\newcommand{\hlg}[1]{\textcolor{green!40!black}{#1}}

\newenvironment{findingbox}[1]
  {\par\medskip\noindent{\color{warnborder}\rule{\linewidth}{0.7pt}}\par
   \vspace{3pt}\noindent\textbf{\color{warnborder}Hallazgo: #1}\par\smallskip\noindent}
  {\par\vspace{3pt}\noindent{\color{warnborder}\rule{\linewidth}{0.7pt}}\medskip}


\title{\textbf{FinHive}\\[4pt]
\large De la teoria a la practica: una arquitectura multiagente jerarquica\\
financiera, construida sobre Databricks}
\author{Matias Fabian Adell}
\date{Agosto 2026}

\begin{document}

\maketitle

\begin{abstract}
\noindent
Este articulo documenta el diseno, la implementacion y la puesta en produccion de
FinHive: un sistema multiagente \textbf{jerarquico} de analisis financiero --- un
supervisor raiz que coordina cinco supervisores de dominio (macroeconomia,
equity research, portfolio \& risk, news \& sentiment, y cripto/activos
alternativos), cada uno con sus propios workers ReAct --- construido sobre
LangGraph y desplegado enteramente sobre Databricks Free Edition. El proyecto
nace como aplicacion practica de una investigacion teorica previa sobre
arquitecturas agenticas (timeline de 24 arquitecturas, 2020--2026) y busca
demostrar, con un sistema real y verificado extremo a extremo, dominio de los
patrones centrales del campo: ReAct, memoria persistente, RAG correctivo,
topologias de orquestacion multiagente, protocolos de interoperabilidad, y
gobernanza de modelos via un LLM Gateway. Mas alla de la arquitectura, este
articulo documenta con el mismo nivel de detalle los \textbf{fallos reales
encontrados y corregidos} durante la construccion --- alucinaciones causadas
por fallos silenciosos de herramientas, bugs de compatibilidad entre
plataformas, errores de ruteo entre dominios --- porque ese proceso de
depuracion es, en la practica, donde mas se pone a prueba el conocimiento de
como estos sistemas realmente fallan.
\end{abstract}

\tableofcontents

%============================================================
\section{Introduccion}
%============================================================

Un modelo de lenguaje individual, por capaz que sea, tiene limites claros
frente a un dominio tan amplio y heterogeneo como el analisis financiero:
requiere datos macroeconomicos estructurados, fundamentals de empresas
cotizantes, calculo cuantitativo de riesgo de portfolio, sentimiento de
mercado en tiempo real, y datos de mercados de criptoactivos --- cada uno con
sus propias fuentes, formatos y ventanas temporales. La respuesta arquitectonica
mas natural a esa heterogeneidad no es un unico agente generalista, sino una
\textbf{red de agentes especializados}, coordinados bajo una topologia que
distribuya la responsabilidad sin perder coherencia en la respuesta final.

FinHive implementa exactamente ese patron: una jerarquia de dos niveles donde
un supervisor raiz decide, por cada turno de conversacion, a cual de cinco
equipos de dominio delegar, y cada equipo de dominio es a su vez un
supervisor que coordina entre dos y tres workers ReAct especializados. El
proyecto se aloja integramente en \textbf{Databricks Free Edition}, aprovechando
Unity Catalog, Foundation Model APIs nativos, Unity AI Gateway, y Databricks
Repos --- con el objetivo explicito de que el costo total de desarrollo y
demostracion sea \$0.

Este articulo tiene tres objetivos. Primero, documentar la arquitectura y las
decisiones de diseno que la sostienen. Segundo, mapear explicitamente que
conceptos de la literatura de arquitecturas agenticas se aplican en que parte
del sistema --- no como referencia decorativa, sino como decisiones de diseno
verificables. Tercero, y quizas mas valioso, documentar los \textbf{problemas
reales} encontrados al construir un sistema de este tipo contra APIs y
modelos reales, no simulados: qué falla, por qué falla, y cómo se corrigió.

%============================================================
\section{Contexto: de la teoria a un caso de uso real}
%============================================================

Este proyecto es la continuacion practica de un trabajo teorico previo
--- un timeline de 24 arquitecturas agenticas organizadas en seis eras
narrativas, desde el RAG clasico (\cite{lewis2020}) hasta los primeros
intentos de sistematizar el campo en una taxonomia formal --- que identifica
una tendencia central: cada era del campo le devuelve control de diseno a un
pipeline fijo y se lo entrega al propio sistema. Primero el modelo gana
agencia sobre su propio control flow (ReAct, \cite{yao2023}); despues esa
agencia se aplica al pipeline de RAG (Self-RAG, CRAG); despues se distribuye
entre multiples agentes especializados bajo topologias de orquestacion
(Network, Supervisor, Jerarquica); y finalmente se estandariza como esa
agencia distribuida interopera entre organizaciones (MCP, A2A).

FinHive es, en ese sentido, un ejercicio deliberado de aplicar esa cadena
completa a un dominio concreto en vez de quedarse en el nivel conceptual: la
pregunta que organiza el proyecto no es \emph{``que arquitectura elegir''}
sino \emph{``que tan poco control de diseno hace falta fijar a mano''} ---
topologia, memoria, tolerancia a fallos --- delegando cada vez mas al propio
sistema, en la misma direccion que senala la investigacion teorica previa.

%============================================================
\section{Arquitectura del sistema}
%============================================================

\subsection{Vision general: supervisor de supervisores}

La Figura~\ref{fig:arquitectura} resume la topologia completa. El patron es
el mismo, replicado en dos niveles: un nodo supervisor decide, con
\emph{structured output}, a que nodo delegar el siguiente turno, hasta decidir
\texttt{FINISH}. En el nivel superior, los ``trabajadores'' son subgrafos
completos (cada supervisor de dominio); en el nivel inferior, son workers
ReAct individuales con tools propias.

\begin{figure}[htb]
\centering
\begin{tikzpicture}[
  node distance=10mm and 6mm,
  every node/.style={font=\small},
  sup/.style={draw=boxborder, fill=boxblue, rounded corners, minimum width=6.6cm,
              minimum height=9mm, align=center, font=\small\bfseries},
  dom/.style={draw=boxborder, fill=white, rounded corners, minimum width=2.55cm,
              minimum height=13mm, align=center, font=\footnotesize},
  wrk/.style={draw=gray, fill=gray!8, rounded corners, minimum width=2.55cm,
              minimum height=6mm, align=center, font=\scriptsize},
  arr/.style={-{Stealth[length=2mm]}, draw=boxborder}
]
\node[sup] (top) {Top-Level Supervisor\\ \footnotesize (routing real via Unity AI Gateway, ADR 0009/0010)};

\node[dom, below left=14mm and 62mm of top] (macro) {Macro};
\node[dom, right=4mm of macro] (equity) {Equity\\Research};
\node[dom, right=4mm of equity] (portfolio) {Portfolio\\\& Risk};
\node[dom, right=4mm of portfolio] (news) {News \&\\Sentiment};
\node[dom, right=4mm of news] (crypto) {Crypto\\\& Alt};

\draw[arr] (top) -- (macro);
\draw[arr] (top) -- (equity);
\draw[arr] (top) -- (portfolio);
\draw[arr] (top) -- (news);
\draw[arr] (top) -- (crypto);

\node[wrk, below=6mm of macro] (m1) {3 workers ReAct};
\node[wrk, below=6mm of equity] (e1) {3 workers ReAct};
\node[wrk, below=6mm of portfolio] (p1) {3 workers ReAct};
\node[wrk, below=6mm of news] (n1) {3 workers ReAct};
\node[wrk, below=6mm of crypto] (c1) {2 workers ReAct};

\draw[arr] (macro) -- (m1);
\draw[arr] (equity) -- (e1);
\draw[arr] (portfolio) -- (p1);
\draw[arr] (news) -- (n1);
\draw[arr] (crypto) -- (c1);
\end{tikzpicture}
\caption{Topologia jerarquica de FinHive: un supervisor raiz enruta entre 5
equipos de dominio (patron ``Supervisor''); cada equipo de dominio enruta
entre sus propios workers ReAct (mismo patron, un nivel abajo). 13 workers en
total sobre 25 tools registradas en Unity Catalog.}
\label{fig:arquitectura}
\end{figure}

\subsection{Los cinco dominios}

\begin{itemize}[leftmargin=*]
  \item \textbf{Macro \& Politica Monetaria}: tasas de interes, inflacion e
  indicadores generales, sobre datos reales de FRED (Federal Reserve Economic
  Data). Workers: \texttt{rates\_worker}, \texttt{inflation\_worker},
  \texttt{indicators\_worker}.
  \item \textbf{Equity Research \& Fundamentals}: cotizaciones, fundamentals
  y filings regulatorios de empresas cotizantes, via yfinance y la API de SEC
  EDGAR (submissions y datos XBRL estructurados). Workers:
  \texttt{fundamentals\_worker}, \texttt{technical\_worker},
  \texttt{filings\_worker}.
  \item \textbf{Portfolio \& Risk Management}: el unico dominio de
  \emph{computo propio} en vez de passthrough a una API --- volatilidad
  anualizada, VaR historico, matriz de correlacion y Sharpe ratio, calculados
  con \texttt{numpy}/\texttt{pandas} sobre precios reales. Workers:
  \texttt{allocation\_worker}, \texttt{risk\_worker}, \texttt{math\_worker}.
  \item \textbf{News \& Sentiment}: noticias con sentimiento estructurado y
  calendario de earnings via Alpha Vantage, con busqueda web (Tavily) como
  \emph{fallback} estilo CRAG cuando la fuente estructurada no alcanza.
  Workers: \texttt{news\_worker}, \texttt{sentiment\_worker},
  \texttt{calendar\_worker}.
  \item \textbf{Crypto \& Alternative Assets}: precio, historial y tendencias
  de criptomonedas via CoinGecko (API publica, sin autenticacion). Workers:
  \texttt{market\_data\_worker}, \texttt{alt\_data\_worker}.
\end{itemize}

\subsection{Infraestructura sobre Databricks}

Todo el sistema corre sobre \textbf{Databricks Free Edition}, con costo total
verificado de \$0:

\begin{itemize}[leftmargin=*]
  \item \textbf{Foundation Model APIs nativos} (\texttt{system.ai.*}) para los
  LLMs de dominio: Llama 3.3 70B (supervisores) y Llama 3.1 8B (workers), sin
  ninguna key externa.
  \item \textbf{Unity Catalog Functions}: las 25 tools de datos estan
  registradas como funciones de Unity Catalog (\texttt{workspace.finhive.*})
  para gobernanza y catalogacion --- el equivalente conceptual a MCP sin
  pasar por los Managed MCP servers de Databricks, que facturan computo
  serverless por invocacion.
  \item \textbf{Unity AI Gateway}: el supervisor raiz --- el nodo mas critico
  del grafo --- corre sobre un \emph{model service} propio
  (\texttt{workspace.finhive.finhive\_router}) con \textbf{routing real}
  entre dos modelos (70\% Llama 3.3 70B / 30\% GPT OSS 120B), mas rate
  limits explicitos sobre los tres endpoints del proyecto.
  \item \textbf{Databricks Repos + notebooks nativos}: el repositorio de
  GitHub esta conectado como Repo en el workspace; un notebook demo
  (\texttt{00\_demo.py}) instala el paquete en modo editable y corre el
  sistema completo dentro de Databricks, con credenciales cargadas via
  \texttt{dbutils.secrets} y \texttt{dbutils.notebook} en vez de archivos de
  configuracion locales.
  \item \textbf{MLflow Tracing}: cada invocacion del grafo genera una traza
  completa (cada nodo, cada tool call, cada llamada a LLM) via
  \texttt{mlflow.langchain.autolog()}.
\end{itemize}

%============================================================
\section{Mapa de conceptos teoricos aplicados}
%============================================================

La Tabla~\ref{tab:conceptos} conecta explicitamente cada concepto de la
investigacion teorica previa con su implementacion concreta en FinHive --- el
objetivo declarado del proyecto no es solo construir un sistema que funcione,
sino demostrar, de forma verificable, el manejo de estos patrones.

\begin{longtable}{@{}p{4.6cm}p{9.7cm}@{}}
\caption{Mapa de trazabilidad teoria--implementacion.}\label{tab:conceptos}\\
\toprule
\textbf{Concepto (papers/seccion)} & \textbf{Donde se aplica en FinHive} \\
\midrule
\endhead
ReAct \cite{yao2023} & Los 13 workers hoja usan
\texttt{langchain.agents.create\_agent} (patron thought--action--observation). \\
Reflexion \cite{shinn2023} & Prompt de auto-critica en los supervisores de
dominio: no aceptar un dato sin que provenga de una tool real. \\
Generative Agents / MemGPT \cite{park2023,packer2023} & Memoria de dos
niveles planificada (checkpointer de conversacion + memoria persistente de
perfil), trabajo futuro documentado en la seccion~\ref{sec:futuro}. \\
Self-RAG / Corrective RAG \cite{asai2024,yan2024} & \texttt{web\_search\_news}
(Tavily) como fallback cuando Alpha Vantage no cubre una consulta ---
patron CRAG de ``no confiar ciegamente en la fuente primaria''. \\
Multi-Agent Network / Supervisor / Jerarquica \cite{wu2023,wang2025} & Los
tres patrones del proyecto de referencia, aplicados literalmente: cada
dominio es un ``Supervisor'' de sus workers; el conjunto completo es
``Jerarquico'' (TalkHier) de dos niveles. \\
Mixture-of-Agents \cite{wangmoa2024} & Sintesis del supervisor raiz cuando
una pregunta cruza dos o mas dominios, combinando las respuestas de equipos
independientes en una unica respuesta coherente. \\
Model Context Protocol (MCP) & Unity Catalog Functions expuestas con
interfaz gobernada y catalogada, en vez de los Managed MCP servers
(decision de costo, ver ADR 0004). \\
LLM Gateway (\S19, resumen conceptual del bootcamp) & Unity AI Gateway:
routing real 70/30 entre dos modelos, rate limits explicitos, usage
tracking --- las mismas responsabilidades (routing/fallback, cost tracking,
policy enforcement) descriptas en la literatura, verificadas en un sistema
real. \\
Adaptive-RAG \cite{jeong2024} & El router del supervisor raiz decide, por
turno, si delegar a un unico equipo o a varios --- version simplificada de
enrutar segun la complejidad real de la pregunta. \\
LLM-as-judge / evaluacion & Trabajo futuro directo: Mosaic AI Agent
Evaluation (framework CLEARS) sobre el grafo ya instrumentado con MLflow. \\
\bottomrule
\end{longtable}

%============================================================
\section{Decisiones de diseno y hallazgos tecnicos reales}
%============================================================

Esta seccion documenta, en orden cronologico, los problemas reales
encontrados construyendo el sistema contra APIs y modelos reales --- no
como anecdota, sino porque en la practica es donde mas se pone a prueba la
comprension de como fallan estos sistemas. El registro completo esta en diez
\emph{Architecture Decision Records} (ADRs) en el repositorio; esta seccion
resume los mas significativos.

\subsection{El costo del LLM: tres decisiones antes de llegar a la correcta}

La primera decision de arquitectura --- que LLM usar --- paso por tres
iteraciones reales antes de asentarse. Primero se evaluo Anthropic via
\emph{External Models} de Databricks; al confirmar que la API de Anthropic
es pay-per-uso separada de una suscripcion de Claude, se paso a Groq (free
tier real, modelos open-weight). Groq quedo configurado y probado, pero al
revisar el catalogo de Foundation Model APIs \emph{nativos} de Databricks se
confirmo, con una llamada real contra el workspace, que estan incluidos sin
costo en Free Edition --- sin necesitar ninguna cuenta de terceros. Se migro
la arquitectura completa a los modelos nativos, eliminando toda dependencia
de key externa para el LLM.

\begin{findingbox}{decisiones de proveedor de LLM no son solo tecnicas}
El costo real de un proveedor de LLM no siempre es evidente hasta probarlo
contra la cuenta real: la documentacion generica no sustituye una llamada de
prueba contra el propio workspace. Las tres iteraciones quedaron documentadas
en ADRs sucesivos en vez de sobrescribir la decision anterior --- el registro
del cambio es en si mismo parte del valor del documento.
\end{findingbox}

\subsection{Un bug de subprocess causa una alucinacion}
\label{sec:hallazgo-windows}

Al conectar las tools de datos como \emph{Unity Catalog Functions}, el modo
recomendado de ejecucion (\texttt{UCFunctionToolkit} en modo \texttt{local})
fallo silenciosamente en Windows: el runner que genera internamente hace
\texttt{import resource}, un modulo de la libreria estandar que \textbf{solo
existe en sistemas Unix}. El fallo se manifesto como un mensaje de error
generico (\emph{``the function execution has been terminated with a
signal''}), y el worker, en vez de admitir que no tenia el dato, \textbf{alucino
una cifra plausible} (5.25\% de tasa de fondos federales, cuando el valor
real --- verificado por separado --- era 3.63\%).

\begin{findingbox}{un fallo silencioso de tool es indistinguible de un dato real}
La correccion tecnica fue simple (envolver las mismas funciones directamente
como tools de LangChain, sin pasar por el sandbox de subprocess), pero el
hallazgo de fondo es mas general: cuando una tool falla sin una senal clara,
el modelo de lenguaje puede optar por una respuesta plausible en vez de una
admision de incertidumbre. Ese comportamiento es exactamente lo que
Corrective RAG y los guardrails de \emph{groundedness} buscan prevenir --- y
esta fue evidencia empirica de por que importan, no solo una referencia
teorica.
\end{findingbox}

\subsection{Un supervisor que no sabia cuando parar}

Al implementar el supervisor raiz con \emph{structured output} para decidir
el proximo equipo, dos problemas de compatibilidad y de disenio de prompt
aparecieron en la misma iteracion. Primero, un \texttt{Literal} construido
dinamicamente con \emph{unpacking} en runtime rompio la conversion de schema
de \texttt{with\_structured\_output} (\texttt{TypeError:\allowbreak\ issubclass()\allowbreak\ arg\ 1\ must\allowbreak\ be\ a\allowbreak\ class}) ---
resuelto usando un campo \texttt{str} simple, validado en codigo en vez de
delegarlo al schema. Segundo, y mas relevante: sobre una pregunta ya
respondida completamente en la primera vuelta, el supervisor re-ruteo al
mismo equipo \textbf{nueve veces} antes de decidir \texttt{FINISH}, gastando
cuota de requests sin necesidad.

La correccion combino dos capas, deliberadamente redundantes: un prompt mas
explicito sobre cuando parar, y un \textbf{limite duro de iteraciones}
independiente del juicio del modelo. Con ambas, la misma clase de pregunta
paso de nueve vueltas a una.

\subsection{Ruteo incorrecto entre dominios y su costo}

Con cuatro dominios reales, la pregunta \emph{``¿cuando es el proximo
earnings de Apple?''} se enruto al equipo de \textbf{Equity} en vez de
\textbf{News \& Sentiment}. El router solo recibia los \emph{nombres} de los
equipos, sin descripcion de sus limites --- nada distinguia ``resultados ya
reportados'' de ``fecha de un resultado futuro''. Equity, sin la tool de
calendario, \textbf{alucino una fecha incorrecta} (2025 en vez de la real,
octubre de 2026) en vez de admitir que no tenia el dato: el mismo patron de
fondo que el hallazgo de la seccion~\ref{sec:hallazgo-windows}, esta vez
causado por una ambiguedad de ruteo en vez de un bug de plataforma.

La correccion fue agregarle al router descripciones explicitas de cada
equipo, incluyendo lineas de ``esto NO es de este equipo'' en los casos de
frontera ya identificados --- un patron que se aplico preventivamente al
sumar el quinto dominio (Crypto), anotando de antemano su frontera con
Equity antes de que apareciera el mismo tipo de bug.

\subsection{Tools defensivas: cuando una API externa falla}

Probando el dominio de Crypto bajo uso intensivo, un \texttt{429 Too Many
Requests} de CoinGecko sin capturar \textbf{crasheo la invocacion completa}
del grafo jerarquico de cinco equipos --- la excepcion se propagaba en vez de
quedar como una observacion que el worker pudiera manejar. El problema era
transversal, no especifico de un dominio: los cinco usan el mismo patron
\texttt{requests.get(...);\ response.raise\_for\_status()}. La correccion fue
un decorador (\texttt{safe\_tool}) que envuelve cualquier tool y convierte
cualquier excepcion en un string de error descriptivo, aplicado en las 25
tools de los cinco dominios.

\subsection{Model routing real, verificado empiricamente}

El hallazgo mas reciente no fue un bug sino una capacidad nueva: Unity AI
Gateway (disponibilidad general, agosto de 2026) permite crear
\emph{model services} de Unity Catalog con routing real entre modelos. Se
configuro \texttt{finhive\_router} con 70\% de trafico a Llama 3.3 70B y
30\% a GPT OSS 120B, y se \textbf{verifico empiricamente}, no solo por
configuracion: seis llamadas de prueba devolvieron cinco respuestas servidas
por Llama y una por GPT OSS 120B, consistente con el split configurado.
El cliente de LangChain especializado de Databricks (\texttt{ChatDatabricks})
todavia no soporta el path nuevo de API que exige este mecanismo; la
integracion se resolvio usando el cliente OpenAI-compatible que Databricks
expone directamente (\texttt{langchain\_openai.ChatOpenAI}, apuntando a
\texttt{/ai-gateway/mlflow/v1}), sin esperar a que la libreria especializada
alcance a la plataforma.

%============================================================
\section{Resultados}
%============================================================

\begin{itemize}[leftmargin=*]
  \item Cinco dominios completos, verificados extremo a extremo contra
  Databricks y APIs de datos reales --- no simulados.
  \item 13 workers ReAct sobre 25 Unity Catalog Functions registradas y
  gobernadas.
  \item Seis tests de integracion pasando juntos, cada uno verificando que
  la respuesta este \emph{grounded} en una llamada real a una tool (no solo
  que el codigo no arroje una excepcion).
  \item Costo total de LLM y embeddings: \textbf{\$0}, verificado en vivo
  contra la cuenta real de Databricks Free Edition, no solo asumido por
  documentacion.
  \item Diez Architecture Decision Records documentando cada decision y cada
  hallazgo con su razonamiento, no solo el estado final.
\end{itemize}

%============================================================
\section{Trabajo futuro}
\label{sec:futuro}
%============================================================

\begin{itemize}[leftmargin=*]
  \item \textbf{Guardrails formales}: tópico, seguridad, \emph{groundedness}
  y anti-jailbreak, con disclaimers explicitos --- motivado directamente por
  los hallazgos de alucinacion documentados en este articulo.
  \item \textbf{Memoria persistente}: checkpointer de conversacion mas
  memoria de largo plazo de perfil/portfolio (patron MemGPT), pendiente
  desde el diseno original.
  \item \textbf{RAG sobre filings completos}: arbol de resumenes al estilo
  RAPTOR sobre el texto completo de 10-K/10-Q, hoy cubierto solo por
  metadata y datos XBRL estructurados.
  \item \textbf{Mosaic AI Agent Framework}: envolver el grafo existente
  (sin reescribirlo) con la interfaz \texttt{ChatAgent} de MLflow y evaluar
  con el framework CLEARS (Correctness, Latency, Execution, Adherence,
  Relevance, Safety) --- extension directa de la instrumentacion de MLflow
  ya activa.
  \item \textbf{Demo desplegada}: aplicacion Streamlit como Databricks App,
  con visualizacion de que equipo/worker respondio cada parte de la
  pregunta.
\end{itemize}

%============================================================
\section{Conclusion}
%============================================================

FinHive demuestra que los patrones centrales de la literatura de
arquitecturas agenticas --- ReAct, RAG correctivo, topologias de
orquestacion jerarquica, protocolos de gobernanza de herramientas y
modelos --- no son solo conceptos de paper, sino decisiones de ingenieria
con consecuencias concretas y verificables cuando se aplican contra sistemas
reales. El valor mas duradero del proyecto no es la arquitectura final en si
misma, sino el registro explicito de cada decision y cada fallo real que
llevo hasta ahi: una alucinacion causada por un bug de plataforma, un loop
sin limite, un error de ruteo entre dominios --- cada uno documentado con su
causa raiz y su correccion, no solo su resultado. Es esa disciplina de
documentacion, mas que cualquier eleccion de framework, la que hace que el
sistema sea auditable y extensible.

%============================================================
\begin{thebibliography}{99}

\bibitem{lewis2020} Lewis, P. et al. (2020). Retrieval-Augmented Generation
for Knowledge-Intensive NLP Tasks. \emph{NeurIPS 33}, 9459--9474.

\bibitem{yao2023} Yao, S. et al. (2023). ReAct: Synergizing Reasoning and
Acting in Language Models. \emph{ICLR}.

\bibitem{shinn2023} Shinn, N. et al. (2023). Reflexion: Language Agents with
Verbal Reinforcement Learning. \emph{NeurIPS}.

\bibitem{park2023} Park, J. S. et al. (2023). Generative Agents: Interactive
Simulacra of Human Behavior. \emph{UIST}.

\bibitem{packer2023} Packer, C. et al. (2023). MemGPT: Towards LLMs as
Operating Systems. \emph{arXiv:2310.08560}.

\bibitem{asai2024} Asai, A. et al. (2024). Self-RAG: Learning to Retrieve,
Generate, and Critique through Self-Reflection. \emph{ICLR}.

\bibitem{yan2024} Yan, S.-Q. et al. (2024). Corrective Retrieval Augmented
Generation. \emph{arXiv:2401.15884}.

\bibitem{jeong2024} Jeong, S. et al. (2024). Adaptive-RAG: Learning to Adapt
Retrieval-Augmented Large Language Models through Question Complexity.
\emph{NAACL}.

\bibitem{wu2023} Wu, Q. et al. (2023). AutoGen: Enabling Next-Gen LLM
Applications via Multi-Agent Conversation. \emph{arXiv:2308.08155}.

\bibitem{wang2025} Wang, Z. et al. (2025). Talk Structurally, Act
Hierarchically: A Collaborative Framework for LLM Multi-Agent Systems.
\emph{arXiv preprint}.

\bibitem{wangmoa2024} Wang, J. et al. (2024). Mixture-of-Agents Enhances
Large Language Model Capabilities. \emph{arXiv:2406.04692}.

\bibitem{adell2026timeline} Adell, M. (2026). Del bibliotecario solitario al
equipo de agentes: una linea de tiempo de las arquitecturas agenticas, de
RAG a los sistemas multiagente (2020--2026). Seminario de Tendencias de
Ciencia de Datos.

\end{thebibliography}

\end{document}
